% 说明：本书的来源、编排与标注，讲义勘误
\frontsection{说明}

这里收的是图论课讲义中不平凡的部分：讲义没有证明或证明不完整的定理，不平凡的例题与课后题的解答，以及对其中经典方法的归纳。定义和简单的题（画图、写矩阵、直接套定义）不收。

内容按方法编排，不按讲义的章节。每一讲讲一种（或两三种相关的）方法：开头说明方法的要点，然后是例题；讲后的作业是新出的练习，用来熟悉这些方法，每题后面都有解答，网页上默认折叠。

\subsection*{标注}

出自讲义的内容，在开头用括号注明，如“（讲义定理 1.4）”“（讲义习题 1.8）”“（讲义问答题 1.3）”。没有注明的是补充的。

\subsection*{讲义勘误}

下面几道讲义习题的原文有误。每题开头的注记给出反例与猜测的原题，解答证明的是猜测的版本。

\begin{enumerate}
\item 习题 1.9（\wegref{we:longest-paths:dirac}）：“长度至少为 $\min\{2\delta,\abs G\}$ 的路或圈”对 $K_2$ 不成立；应为“长度至少为 $\min\{2\delta,\abs G-1\}$ 的路”。
\item 习题 1.12（\wegref{we:double-counting:games}）：“$n$ 个球队赛完 $n-1$ 场，必有一队至少赛过 $3$ 场”不成立（各队排成一列，相邻两队各赛一场）；应为“赛完 $n+1$ 场”。
\item 习题 1.17（\wegref{we:distance-layers:equal-distance}）：对不连通的图不成立（$K_2\cup K_1$）；应限定 $G$ 连通。
\end{enumerate}
另有一处措辞需要澄清：习题 1.8（\wegref{we:longest-paths:sqrt-k}）中“连接它们之间的路的长度至少为 $k$”应理解为存在这样一条路，而不是两点的距离至少为 $k$。

\subsection*{没有收的讲义习题}

第 1 章：习题 1、2、3、5、6、7、14、15、19、20，问答题 1.1、1.2、1--4、1--5。
